\UTIResumenPagina
  {spanish}
  {Resumen Ejecutivo}
  {La Ley Orgánica de Protección de Datos Personales obliga a todo responsable a utilizar una metodología de análisis y gestión de riesgos, y su omisión constituye una infracción grave. En Ecuador, 16\,161 instituciones educativas y 280 instituciones de educación superior están sujetas a esa obligación, mientras que los delegados de protección de datos registrados en abril de 2025 equivalían como máximo al 7,0\,\% de ellas y no existe una medición oficial del cumplimiento de las universidades. El objetivo general fue diseñar una metodología de diagnóstico y gestión de riesgos para el tratamiento de datos personales de estudiantes y docentes en instituciones de educación superior, integrando la guía de gestión de riesgos de la Superintendencia de Protección de Datos Personales con la norma ISO/IEC 27005:2022. Se siguió la investigación en ciencia del diseño, con enfoque mixto, aplicada al proceso de admisión y matrícula de la Universidad Tecnológica Indoamérica, sede Quito. La metodología comprende una fase de preparación, una de diagnóstico de madurez y cuatro de gestión (identificación, evaluación, mitigación y monitoreo), equivalentes a ISO/IEC 27035-1:2023 y NIST SP 800-61r3; estima la probabilidad anual y el impacto sobre el titular mediante un panel Delphi calibrado con el puntaje de Brier, y asigna el nivel de riesgo con una tabla de decisión que impide ocultar daños graves tras probabilidades bajas. Se determinaron diez requisitos para ejecutar el diagnóstico, trece escenarios candidatos trazables a la ley y una plantilla de matriz cuyas 1\,057 fórmulas se evalúan sin errores y que supera seis casos de prueba. Los resultados de la aplicación en la institución y de la validación por expertos dependen de la autorización institucional y quedan identificados como pendientes. Se concluye que la metodología está completamente especificada y que su validez se juzgará con doce indicadores con metas definidos de antemano.}
  {Descriptores}
  {evaluación de impacto, gestión de riesgos, instituciones de educación superior, protección de datos personales}
